% Generated draft: requires mathematical and source review.
\documentclass[11pt]{article}
\usepackage[margin=25mm]{geometry}
\usepackage{fontspec}
\setmainfont{Latin Modern Roman}
\usepackage{amsmath,amssymb,mathtools,unicode-math}
\setmathfont{Latin Modern Math}
\usepackage{enumitem,needspace,xurl,hyperref}
\setlength{\parindent}{0pt}
\setlength{\parskip}{5pt}
\setlength{\emergencystretch}{3em}
\newcommand{\N}{\mathbb{N}}
\newcommand{\Z}{\mathbb{Z}}
\newcommand{\Q}{\mathbb{Q}}
\newcommand{\R}{\mathbb{R}}
\newcommand{\C}{\mathbb{C}}
\newcommand{\F}{\mathbb{F}}
\newcommand{\A}{\mathbb{A}}
\newcommand{\PP}{\mathbb{P}}
\newcommand{\E}{\mathbb{E}}
\newcommand{\eps}{\varepsilon}
\newcommand{\dd}{\,\mathrm{d}}
\begin{document}
\section*{UnsolvedMath ID: 700667}
\textbf{Source problem code:} catalogue entry
\subsection{Definitions and mathematical statement}
Let \(p\) be a prime. For an abelian group \(A\), written additively, the \(p\)-adic topology has a neighbourhood basis at \(0\) consisting of the subgroups
\[
p^nA=\{p^na:a\in A\},\qquad n\geq 1.
\]
Its separated completion is
\[
\widehat A_p=\varprojlim_{n\geq 1} A/p^nA,
\]
with each quotient given the discrete topology. Write \(\tau_{\mathrm{comp}}\) for the resulting inverse-limit topology on \(\widehat A_p\), whose basic neighbourhoods of \(0\) are the kernels of the coordinate projections. Write \(\tau_p(\widehat A_p)\) for the intrinsic \(p\)-adic topology on its underlying abelian group, defined using the subgroups \(p^n\widehat A_p\). These definitions do not require the original group \(A\) to be separated in its \(p\)-adic topology.

The supplied source reports the abelian-group assertion
\[
\tau_{\mathrm{comp}}=\tau_p(\widehat A_p),
\]
and attributes the analogous assertion for nilpotent groups to R. Warfield. These are source-reported assertions; no proof or independent verification is supplied here.

\textbf{Problem.} Find a categorical setting for the assertion that the topology on the completion of a group in its \(p\)-adic topology is itself the \(p\)-adic topology, with a scope that includes nilpotent groups. The proposed setting must specify the relevant category, the meaning of the \(p\)-adic topology and completion, and the hypotheses under which the two topologies agree.

\textbf{Unresolved qualifications.} The source does not define the nonabelian \(p\)-adic topology or specify the hypotheses of Warfield's result. It also does not prescribe the categorical structure or the precise sense in which the category of nilpotent groups must be included. Consequently, the categorical request is a research target rather than a fully specified quantified proposition; these missing specifications must be resolved before a precise categorical theorem can be stated.
\subsection{Short English statement}
Find a categorical framework explaining the source-reported fact that completing an abelian group in its \(p\)-adic topology produces a group whose completion topology is again its intrinsic \(p\)-adic topology, and covering the analogous result attributed to R. Warfield for nilpotent groups. The nonabelian topology, hypotheses, and intended categorical framework are not specified in the supplied record.
\subsection{Sources}
\begin{itemize}
\item Supplied local UnsolvedMath registry record: permanent ID \(700667\), exact source problem code \texttt{catalogue entry}, title \emph{Kourovka Notebook Problem 10.52}.
\item \emph{The Kourovka Notebook (Unsolved Problems in Group Theory)}, original problem 10.52, 10th issue (1986); supplied reference to the 21st issue edition (2026), arXiv:1401.0300v48. The registry describes the collection as edited by E. I. Khukhro and V. D. Mazurov.
\item Attribution as supplied: the registry and background identify S. V. Rychkov as proposer, with proposed year 1986, whereas the statement begins ``(R. Mines).'' The nilpotent-group result is attributed in the statement to R. Warfield, without a separate bibliographic citation. The respective attribution roles are unresolved.
\item Supplied discussion: an explicitly unverified AI solution claim by Achyuth Jayadevan, \emph{Preprint} (2026), \url{https://zenodo.org/records/22852281}. The supplied description says that a categorical setting is found; this claim is not treated here as an established solution.
\item Supplied status review dated exactly \texttt{2026-10-16}: canonical status \texttt{open}, evidence label \texttt{unrefereed\_claim}. Its notes report a \(\#\) marker for an unverified AI solution claim in the October 2026 21st issue, arXiv:1401.0300v48. This assessment is supplied dated curation, not independently verified source or literature checking.
\end{itemize}
\end{document}
